\chapter{การออกแบบและพัฒนาแอปพลิเคชันบนสมาร์ทโฟน}

\section{สถาปัตยกรรมซอฟต์แวร์แอปพลิเคชัน (Clean Architecture in Flutter)}
แอปพลิเคชัน DurianLeaf AI พัฒนาขึ้นด้วยเฟรมเวิร์ก Flutter โดยใช้ภาษา Dart ภายใต้แบบแผนสถาปัตยกรรม Clean Architecture และการจัดการสถานะด้วย Provider ซึ่งแยกโครงสร้างการทำงานออกเป็น 3 ชั้นหลักอย่างชัดเจน เพื่อให้ระบบมีความเสถียร สามารถทดสอบแยกส่วนได้ (Testability) และรองรับการขยายตัวในอนาคต

\begin{enumerate}
  \item \textbf{Presentation Layer} ประกอบด้วยวิดเจ็ตหน้าจอหลัก (\texttt{HomeScreen}), ตัวแสดงผลสตรีมกล้อง (\texttt{CameraView}), แถบควบคุมพื้นที่ตรวจจับ (\texttt{RoiToolbar}) และแผ่นหน้าต่างเกณฑ์วิชาการ (\texttt{AcademicLegendModal})
  \item \textbf{Domain Layer} บรรจุโมเดลข้อมูลทางพฤกษศาสตร์และพยาธิวิทยา เช่น \texttt{DurianDiseaseClass}, \texttt{NutrientColorStandard}, \texttt{ColorimetryEngine} และตัวคำนวณดัชนีพืชพรรณ
  \item \textbf{Data Layer} จัดการการเชื่อมต่ออุปกรณ์ฮาร์ดแวร์จริง เช่น กล้องถ่ายภาพ (\texttt{CameraPlugin}), เซนเซอร์จีพีเอส (\texttt{Geolocator}), และตัวขับเคลื่อนโมเดลปัญญาประดิษฐ์ (\texttt{TfliteNativeBridge})
\end{enumerate}

\section{ระบบสตรีมภาพสดและท่อประมวลผลความเร็วสูง (Zero-Lag Camera Pipeline)}
การประมวลผลโครงข่ายประสาทเทียมและอัลกอริทึมสีควบคู่ไปกับการเรนเดอร์ภาพสด $30\text{ FPS}$ ในสภาพแปลงเกษตรกลางแจ้ง มักก่อให้เกิดปัญหาความร้อนสะสมและการหน่วงของภาพ (Frame Dropping) ระบบ DurianLeaf AI จึงใช้กลไก Mutex Lock และเฟรมบัฟเฟอร์แบบอะซิงโครนัส

\begin{equation}
  \Delta t_{\text{frame}} \ge \frac{1000\text{ ms}}{\text{Target FPS}} \approx 33.3\text{ ms}
\end{equation}
หากโมเดลปัญญาประดิษฐ์ยังประมวลผลเฟรมเดิมไม่เสร็จ ระบบจะข้ามเฟรมแทรกซ้อนทันทีโดยไม่เข้าคิวสะสม ทำให้การแสดงผลหน้าจอตอบสนองลื่นไหลแบบเรียลไทม์เสมอ

\section{องค์ประกอบของแถบเมนูควบคุมพื้นที่ตรวจจับ (Interactive ROI Floating Toolbar)}
ที่ขอบขวาของหน้าจอหลัก มีแถบลอยสำหรับการควบคุมพื้นที่ตรวจจับแบบอินเตอร์แอคทีฟ (Interactive ROI Floating Toolbar) ซึ่งได้รับการออกแบบตามมาตรฐานเดียวกับเครื่องมือตรวจวัดทางวิชาการ โดยประกอบด้วย 4 ฟังก์ชันหลัก

\subsection{การปรับเปลี่ยนรูปทรงและขนาดของกรอบเล็งเป้า}
\begin{itemize}
  \item \textbf{กรอบสี่เหลี่ยม (Rectangle ROI)} เหมาะสำหรับการกำหนดพื้นที่รอยโรคขนาดใหญ่ เช่น แผลใบไหม้ไฟทอปธอร่า
  \item \textbf{กรอบวงกลม (Circular ROI)} เหมาะสำหรับการสุ่มเก็บข้อมูลเนื้อเยื่อใบสมบูรณ์เพื่อวิเคราะห์ค่า SPAD และธาตุอาหาร
  \item \textbf{ปุ่มขยายและย่อขนาด (+ / -)} เกษตรกรสามารถปรับขนาดเส้นผ่านศูนย์กลางของกรอบเล็งได้ตั้งแต่ $80$ ถึง $340$ พิกเซล เพื่อให้ครอบคลุมแผ่นใบในระยะซูมต่างๆ
\end{itemize}

\subsection{แถบแสดงเฉดสีตามเกณฑ์วิชาการ (Academic Color Strip)}
ใต้ปุ่มควบคุมขนาด จะเป็นแถบเฉดสีแนวตั้งที่แสดงเกณฑ์มาตรฐานทางวิชาการ เกษตรกรสามารถแตะปุ่มสลับโหมดเพื่อเปลี่ยนเกณฑ์สีที่ต้องการมอนิเตอร์สด ได้แก่
\begin{enumerate}
  \item \textbf{โหมดคลอโรฟิลล์ (SPAD Strip)} แสดงแถบไล่เฉด 5 ระดับ ตั้งแต่สีซีดขาว ($<25$) ไปจนถึงสีเขียวมรกตเข้ม ($>55$)
  \item \textbf{โหมดไนโตรเจน (N Strip)} แถบสีแสดงเกณฑ์โภชนาการตั้งแต่เหลืองซีด ($<1.8\%$) จนถึงเขียวเข้มจัด ($>3.2\%$)
  \item \textbf{โหมดโพแทสเซียม (K Strip)} แถบสีส้มอิฐถึงน้ำตาลทองเตือนขอบใบไหม้
  \item \textbf{โหมดความรุนแรงของโรค (DSI Strip)} ไล่เฉดสีเขียว (Grade 0) สีเหลือง (Grade 1-2) สีส้ม (Grade 3) และสีแดงเข้ม (Grade 4)
\end{enumerate}

\subsection{ไฟเรืองแสงนีออนและป้ายชี้เป้าสด (Neon Glow Live Pointer Badge)}
เมื่อกล้องถ่ายภาพจับสัญญาณภาพใบไม้และ AI คำนวณค่าได้ ช่องสีบนแถบแนวตั้งที่ตรงกับระดับผลลัพธ์ปัจจุบันจะเปล่งแสงนีออนเรืองแสงกระพริบ และมีป้ายชี้เป้าสดสีดำขอบทอง เช่น \texttt{$\blacktriangleleft$ SPAD 52.4 (สมบูรณ์)} เลื่อนไปยังตำแหน่งระดับสีนั้นอย่างแม่นยำ

\subsection{หน้าต่างแสดงคำอธิบายเกณฑ์วิชาการ (Academic Legend Modal)}
เมื่อแตะที่แถบเฉดสี ระบบจะเปิดหน้าต่างป๊อปอัปแสดงตารางวิชาการ บรรจุข้อมูลพิกัดสีจริง ($RGB, HEX, CIE\ L^*a^*b^*$) นิยามทางพฤกษศาสตร์ และแนวทางการปฏิบัติดูแลรักษาสำหรับเกษตรกร

\section{ระบบการสอนโมเดลเพิ่มเติมในแปลงจริง (Few-Shot Model Trainer)}
หากเกษตรกรพบรอยโรคแปลกใหม่ในแปลง สามารถแตะไอคอนรูปหมวกปริญญา เพื่อเปิดหน้าต่างสอนโมเดลเพิ่มเติม
\begin{enumerate}
  \item ป้อนชื่อโรคหรืออาการเฉพาะถิ่น
  \item ระบุข้อแนะนำในการกำจัดหรือสูตรปุ๋ย
  \item ส่องกล้องแล้วกดถ่ายภาพตัวอย่าง $2 - 3$ มุมมอง ระบบจะสกัดเวกเตอร์เอกลักษณ์ 128 มิติ
  \item กดปุ่มบันทึก ระบบจะจดจำและตรวจจับอาการดังกล่าวในกล้องสดได้ทันทีโดยไม่ต้องต่ออินเทอร์เน็ต
\end{enumerate}

\section{การบูรณาการสถานีตรวจวัดสภาพแวดล้อม HandySense และดัชนี VPD}
ความสมบูรณ์และการดูดซึมธาตุอาหารของใบไม้ทุเรียนมีความสัมพันธ์อย่างแนบแน่นกับจุลภูมิอากาศระดับแปลง (Microclimate) ระบบ DurianLeaf AI จึงได้เชื่อมโยงเข้ากับสถานีตรวจวัด HandySense และเซนเซอร์อุตสาหกรรมเกษตรแม่นยำ
\begin{enumerate}
  \item \textbf{เซนเซอร์ตรวจวัดบรรยากาศ SHT45} ตรวจวัดอุณหภูมิอากาศ ($T \text{ ในหน่วย } ^\circ\text{C}$) และความชื้นสัมพัทธ์ ($\text{RH } \%$)
  \item \textbf{เซนเซอร์วัดคุณสมบัติดิน Soil 7-in-1} ตรวจวัดความชื้นในดิน ($\% \text{ VWC}$), การนำไฟฟ้า (EC), ความเป็นกรด-ด่าง (pH), และธาตุอาหารดิน N-P-K
  \item \textbf{เซนเซอร์วัดความเข้มแสงแดด BH1750} ตรวจวัดค่าความสว่างแสง ($0 - 65,000 \text{ Lux}$) เพื่อประเมินสภาวะการสังเคราะห์แสง
\end{enumerate}

\subsection{การคำนวณความดันไอน้ำในอากาศส่วนขาดตามสมการ Tetens (Tetens VPD Equation)}
ดัชนี Vapor Pressure Deficit (VPD) เป็นตัวชี้วัดสำคัญที่สุดของแรงดึงการคายน้ำและการเปิด-ปิดปากใบของทุเรียน โดยระบบทำการคำนวณแบบเรียลไทม์จากสมการ Tetens Equation
\begin{align}
  e_s(T) &= 0.61078 \exp\left( \frac{17.27 \times T}{T + 237.3} \right) \quad (\text{kPa}) \\
  e_a(T, \text{RH}) &= e_s(T) \times \frac{\text{RH}}{100} \quad (\text{kPa}) \\
  \text{VPD} &= e_s(T) - e_a(T, \text{RH}) \quad (\text{kPa})
\end{align}
โดยเกณฑ์การประเมินสภาวะทางสรีรวิทยาของทุเรียนแบ่งออกเป็น 4 ระดับ
\begin{itemize}
  \item $\text{VPD} < 0.4\text{ kPa}$ อากาศชื้นจัด การคายน้ำหยุดชะงัก เสี่ยงต่อการระบาดของโรคราใบติดและไฟทอปธอร่า
  \item $0.4 \le \text{VPD} \le 0.8\text{ kPa}$ ความชื้นค่อนข้างสูง พืชคายน้ำปานกลาง เหมาะสำหรับระยะฟื้นต้นทำใบ
  \item $0.8 < \text{VPD} \le 1.2\text{ kPa}$ \textbf{ช่วงสภาวะสมบูรณ์สูงสุด (Optimal Zone)} ปากใบเปิดกว้าง เกิดการดึงน้ำและธาตุอาหารจากรากขึ้นสู่ยอดได้อย่างมีประสิทธิภาพสูงสุด
  \item $\text{VPD} > 1.6\text{ kPa}$ อากาศแห้งแล้งจัด พืชเกิดความเครียด (Water Stress) ปากใบจะปิดตัวลงเพื่อรักษาความชื้น ทำให้การดูดธาตุอาหารหยุดชะงัก ระบบจะส่งสัญญาณเตือนเตือนภัยสีส้มทันที
\end{itemize}

\section{ระบบสมุดบันทึกประวัติสุขภาพรายต้นผูกพิกัดภูมิศาสตร์ (Tree Passport System)}
เพื่อรองรับการทำเกษตรแม่นยำรายต้น (Precision Tree-by-Tree Management) ระบบได้ออกแบบสถาปัตยกรรมข้อมูลบัตรประจำตัวสุขภาพต้นไม้ (\texttt{Tree Passport})
\begin{enumerate}
  \item \textbf{การระบุรหัสต้นไม้ (Tree ID)} เกษตรกรสามารถกำหนดรหัสต้นไม้ เช่น \texttt{TR-DUR-01}, \texttt{TR-DUR-02} เพื่อจำแนกต้นทุเรียนแต่ละต้นในแปลง
  \item \textbf{การบันทึกพิกัดภูมิศาสตร์ระดับความแม่นยำสูง (GPS Geotagging)} ผูกโยงพิกัดละติจูด-ลองจิจูดในระบบพิกัด WGS84 และระดับความสูงจากระดับน้ำทะเลปานกลาง (MSL Altitude)
  \item \textbf{บัตรตรวจสุขภาพรายต้น (Tree Passport Dialog)} สามารถเปิดดูประวัติย้อนหลังของต้นนั้นๆ ได้ทันที แสดงภาพถ่ายตัวอย่าง, ระดับความรุนแรงของโรค, ค่าคลอโรฟิลล์ SPAD, ความเข้มข้นธาตุอาหาร NPK, และสภาพแวดล้อม HandySense ณ วินาทีที่บันทึก
\end{enumerate}

\section{สถาปัตยกรรมการจัดวางหน้าจอที่ไร้การซ้อนทับ (Non-Overlapping Ergonomic Layout)}
จากการทดสอบใช้งานจริงในสภาพแปลง พบว่าการแสดงผลข้อมูลทางวิทยาศาสตร์ที่หลากหลายพร้อมกันบนหน้าจอสมาร์ทโฟนอาจก่อให้เกิดปัญหาชิ้นส่วน UI บดบังกัน จึงได้ออกแบบเลย์เอาต์ตามหลักการยศาสตร์ดิจิทัล (Digital Ergonomics)
\begin{itemize}
  \item \textbf{แถบโทรมาตรบริบทแปลงแบบรวมศูนย์ (Unified AgriPhysics Bar)} รวมข้อมูลรหัสต้นไม้, รุ่นอายุใบ, ระยะพัฒนาการของทรงพุ่ม, และสถานะ HandySense เข้าไว้ในแถบแนวนอนเลื่อนได้แถวเดียว (Single Horizontal Strip) ช่วยประหยัดพื้นที่แนวตั้งได้ถึง $34$ พิกเซล
  \item \textbf{แถบควบคุม ROI แบบกะทัดรัด (Compact Floating ROI Toolbar)} ปรับความกว้างแถบเหลือ $50$ พิกเซล และจัดตำแหน่งให้เริ่มต้นอยู่ใต้แถบข้อมูลด้านบนอย่างสมบูรณ์ ($\text{top} = \text{topPadding} + 144$) ทำให้ไม่บดบังเมนูหัวเรื่องและมีระยะห่างโปร่งสบายจากกล่อง HUD ด้านล่าง
\end{itemize}

\begin{theorybox}
\textbf{การยศาสตร์การมองเห็นกลางแจ้ง (Outdoor Vision Ergonomics)}\par
การใช้งานสมาร์ทโฟนในสวนทุเรียนที่มีแสงแดดจ้าทำให้สายตาแยกแยะความเปรียบต่างได้ยาก การออกแบบ UI จึงใช้ชุดสี Cyber Dark (\texttt{\#0A0F1D}) เป็นพื้นหลัง และใช้เส้นนำสายตาสีนีออนเรืองแสงที่มีความสว่างสูงกว่า $85\%$ ร่วมกับตัวเลขอักษรขนาดใหญ่พิเศษ ทำให้เกษตรกรสามารถอ่านผลลัพธ์ได้อย่างชัดเจนแม้อยู่ใต้แสงแดดจัด
\end{theorybox}

\begin{promptbox}
\textbf{พรอมพ์วิศวกรรมสำหรับการพัฒนาแอปพลิเคชันมือถือและ UI/UX (Prompt Eng-05)}\par
\texttt{"Act as a Senior Flutter and Mobile UI/UX Engineer. Develop a high-performance cross-platform application for real-time durian leaf disease diagnosis, nutrient colorimetry, and microclimate IoT telemetry. Implement Clean Architecture with Provider state management. Integrate Tetens equation for live Vapor Pressure Deficit (VPD) monitoring from HandySense stations. Create a unified non-overlapping ergonomic layout with horizontal telemetry scroll strips and a compact cyber-dark ROI toolbar. Ensure offline SQLite persistence for Tree Passport records, GPS coordinates, and exportable CSV audit logs."}
\end{promptbox}
